% thermo_piston: work done by a gas on a moving piston (Vol II Ch 4, fig:thermo-piston)
% Gas at pressure P in a cylinder closed by a piston of face area A.  The gas pushes
% the piston with force F = PA (crimson, a force).  When the piston moves out by a small
% distance ds (ghost = later position), the gas does work dW = F ds = PA ds = P dV,
% where dV = A ds is the thin slice swept out (the slice style).
\begin{tikzpicture}[x=1cm,y=1cm]
  \def\L{3.4}   % piston face position
  \def\ds{0.8} % displacement
  \def\h{1.5}   % half-height of bore
  % gas
  \fill[cMoon] (0,-\h) rectangle (\L,\h);
  \foreach \x/\y in {0.4/0.8,1.1/-0.6,1.8/0.3,2.6/-1.0,3.1/0.9,0.7/-1.1,2.2/1.1,1.4/0.9,2.9/-0.2,0.5/0.1}
    \node[point mass,cSteel] at (\x,\y) {};
  \node[anchor=center] at (1.75,-0.25) {gas: $P$, $V$};
  % swept slice dV = A ds
  \fill[cAcc!40] (\L,-\h) rectangle (\L+\ds,\h);
  \draw[cAcc,thin line] (\L+\ds,-\h) -- (\L+\ds,\h);
  % cylinder walls (closed left end)
  \draw[outline] (\L+2.3,\h) -- (0,\h) -- (0,-\h) -- (\L+2.3,-\h);
  % ghost of piston at later position
  \begin{scope}[ghost]
    \draw[body] (\L+\ds,-\h) rectangle (\L+\ds+0.16,\h);
  \end{scope}
  % piston (now) and rod
  \draw[body] (\L,-\h) rectangle (\L+0.16,\h);
  \draw[outline,fill=cRock] (\L+0.16,-0.09) rectangle (\L+2.6,0.09);
  % face area label
  \draw[leader] (\L+0.08,-1.05) -- (\L+1.0,-2.0) node[anchor=north west,inner sep=1pt] {piston face, area $A$};
  % force of gas on piston
  \draw[force] (\L+0.16,0.62) -- ++(1.35,0) node[anchor=west] {$F=PA$};
  % displacement
  \draw[dim] (\L,\h+0.32) -- (\L+\ds,\h+0.32);
  \node[cSteel,anchor=south] at (\L+0.5*\ds,\h+0.36) {$ds$};
  \draw[dimtick] (\L,\h) -- (\L,\h+0.45);
  \draw[dimtick] (\L+\ds,\h) -- (\L+\ds,\h+0.45);
  \draw[leader,cAcc] (\L+0.5*\ds,-1.25) -- (\L-0.9,-2.0) node[anchor=north east,inner sep=1pt,cAcc] {slice $dV=A\,ds$};
\end{tikzpicture}
